%==============================================================================
\section{Khái niệm cơ bản}
\label{sec:khai-niem-nen}
%==============================================================================

\subsection{Machine Learning là gì?}
\label{sec:ml-la-gi}

\begin{dinhnghia}[title={Machine Learning (ML)}]
\textbf{Học máy} là một nhánh của \textbf{Trí tuệ nhân tạo (AI)} tập trung vào việc xây dựng thuật toán/mô hình có thể \textbf{học từ dữ liệu} để dự đoán hoặc ra quyết định mà không cần viết quy tắc tường minh cho mọi trường hợp.
\end{dinhnghia}

\begin{ynghia}[title={ML giúp giải bài toán nào?}]
ML đặc biệt hữu ích khi:
\begin{itemize}
  \item Quy tắc quá phức tạp/khó mô tả bằng tay (nhận dạng ảnh/giọng nói, NLP).
  \item Môi trường thay đổi theo thời gian (hành vi người dùng, mạng, rủi ro gian lận).
  \item Cần ra quyết định dựa trên dữ liệu với quy mô lớn (tự động hóa).
\end{itemize}
\end{ynghia}

\subsection{Các thành phần cốt lõi: dữ liệu, đặc trưng, mô hình}
\label{sec:thanh-phan-co-lo}

\begin{dinhnghia}[title={Dữ liệu (data)}]
Dữ liệu là ``nguyên liệu'' của ML. Dữ liệu có thể là có cấu trúc (bảng), bán cấu trúc, hoặc phi cấu trúc (văn bản, ảnh, âm thanh).
Chất lượng và số lượng dữ liệu ảnh hưởng mạnh đến chất lượng mô hình.
\end{dinhnghia}

\begin{dinhnghia}[title={Đặc trưng (feature)}]
\textbf{Feature} là các biến/thuộc tính mô tả đầu vào. Chọn feature phù hợp (feature selection) và tạo feature mới (feature engineering) thường quyết định phần lớn hiệu quả mô hình.
\end{dinhnghia}

\begin{dinhnghia}[title={Mô hình (model)}]
Mô hình là biểu diễn toán học học được mối liên hệ giữa feature và đầu ra (dự đoán/quyết định).
Mô hình được học từ tập train và đánh giá trên tập validation/test để kiểm tra khả năng tổng quát.
\end{dinhnghia}

\subsection{Train / validation / test và overfitting}
\label{sec:train-val-test}

\begin{dinhnghia}[title={Huấn luyện (training)}]
Training là quá trình cho thuật toán xem nhiều ví dụ dữ liệu và điều chỉnh tham số nội bộ để giảm sai số dự đoán.
\end{dinhnghia}

\begin{dinhnghia}[title={Kiểm thử (testing)}]
Testing là đánh giá mô hình trên dữ liệu \textbf{chưa từng thấy} nhằm ước lượng khả năng tổng quát hóa.
\end{dinhnghia}

\begin{dinhnghia}[title={Overfitting và Underfitting}]
\begin{itemize}
  \item \textbf{Overfitting}: mô hình quá phức tạp, bám sát train, kém trên dữ liệu mới.
  \item \textbf{Underfitting}: mô hình quá đơn giản, không học được cấu trúc, kém cả train và test.
\end{itemize}
\end{dinhnghia}

\subsection{Các kiểu học máy}
\label{sec:bon-kieu-hoc}

\begin{dinhnghia}[title={Học có giám sát (Supervised)}]
Dữ liệu có \textbf{nhãn} (đầu ra đúng). Hai nhóm bài toán phổ biến:
\textbf{hồi quy} (đầu ra số) và \textbf{phân loại} (đầu ra lớp).
\end{dinhnghia}

\begin{vidu}[title={Minh hoạ supervised}]
\tsimg{00_khai_niem_nen/img05.png}
\end{vidu}

\begin{dinhnghia}[title={Học không giám sát (Unsupervised)}]
Dữ liệu \textbf{không có nhãn}. target thường là tìm cấu trúc/nhóm trong dữ liệu:
phân cụm, phát hiện bất thường, giảm chiều, \ldots
\end{dinhnghia}

\begin{vidu}[title={Minh hoạ unsupervised}]
\tsimg{00_khai_niem_nen/img06.png}
\end{vidu}

\begin{dinhnghia}[title={Bán giám sát (Semi-supervised)}]
Kết hợp một phần nhỏ dữ liệu có nhãn và phần lớn dữ liệu không nhãn (thường do gán nhãn tốn kém).
\end{dinhnghia}

\begin{vidu}[title={Minh hoạ semi-supervised}]
\tsimg{00_khai_niem_nen/img07.png}
\end{vidu}

\begin{dinhnghia}[title={Tăng cường (Reinforcement)}]
Tác tử (agent) tương tác với môi trường, nhận thưởng/phạt, học chính sách hành động để tối đa hóa phần thưởng dài hạn.
\end{dinhnghia}

\begin{vidu}[title={Minh hoạ reinforcement}]
\tsimg{00_khai_niem_nen/img08.png}
\end{vidu}

\subsection{Chu trình giải bài toán ML}
\label{sec:chu-trinh-ml}

\begin{phuongphap}[title={Chuỗi bước điển hình (nhìn từ góc ``thực hành'')}]
\begin{enumerate}
  \item Xác định bài toán và tiêu chí thành công.
  \item Thu thập và tiền xử lý dữ liệu (làm sạch, thiếu dữ liệu, nhiễu).
  \item Khám phá dữ liệu và trực quan hóa để hiểu phân phối/pattern.
  \item Chọn/thiết kế feature.
  \item Chọn mô hình, huấn luyện, đánh giá.
  \item Tinh chỉnh siêu tham số, kiểm tra lại.
  \item Triển khai, giám sát, cập nhật theo thời gian.
\end{enumerate}
\end{phuongphap}

\begin{vidu}[title={Hình minh hoạ khái niệm nền}]
\tsimg{00_khai_niem_nen/img01.jpg}
\tsimg{00_khai_niem_nen/img03.png}
\tsimg{00_khai_niem_nen/img04.png}
\end{vidu}

